% !TEX root = task-4-5-5.tex
% Editable task bank. Append new BankTask blocks at the end of the document.
\documentclass[12pt,a4paper]{article}
\usepackage{physics-textbook}
\usepackage{graphicx,needspace}
\usetikzlibrary{decorations.pathmorphing}
\newif\ifShowSolutions
\ifdefined\TutorEdition
  \ShowSolutionstrue
\else
  \ShowSolutionsfalse % Change to \ShowSolutionstrue to display solutions.
\fi
\newcounter{banktask}
\newsavebox{\BankTaskBox}
\newcommand{\BankTaskContent}[3]{%
  \noindent\textbf{\thebanktask.}\quad #2\par
  \ifShowSolutions
    \Needspace{50mm}% Keep the solution heading with its opening text or diagram.
    #3\par
    \ifdefined\TutorEdition
      \ifcsname TutorSource#1\endcsname
        \nopagebreak\smallskip{\small\csname TutorSource#1\endcsname\par}%
      \fi
    \fi
  \fi
}
\newcommand{\BankTask}[3]{%
  \refstepcounter{banktask}\par\medskip
  \sbox{\BankTaskBox}{\begin{minipage}{\linewidth}%
    \BankTaskContent{#1}{#2}{#3}\end{minipage}}%
  \ifdim\dimexpr\ht\BankTaskBox+\dp\BankTaskBox\relax>\dimexpr\textheight-20mm\relax
    \Needspace{35mm}\BankTaskContent{#1}{#2}{#3}%
  \else
    \Needspace{\dimexpr\ht\BankTaskBox+\dp\BankTaskBox\relax}%
    \noindent\usebox{\BankTaskBox}\par
  \fi
  \medskip
}
\begin{document}
\BookSubsection{4.5.5}{Ход лучей в рассеивающей линзе}


% Source 29760: https://phys-ege.sdamgia.ru/problem?id=29760
\BankTask{29760}{
Рассеивающая линза с фокусным расстоянием $F_1=10$ см расположена перед собирающей линзой с фокусными расстоянием $F_2=15$ см. Лучи, идущие от точечного источника света, расположенного на расстоянии $d=10$ см от pacceивающей линзы, пройдя систему образовали пучок лучей, параллельный главной оптической оси. Найдите расстояние между линзами.
}{
\smallskip\noindent\textbf{Решение}\par
$1$.  Найдем расстояние от рассеивающей линзы до изображения точечного источника света по формуле тонкой линзы 

\[\frac{1}{d} - \frac{1}{f} = - \frac{1}{F_1} ,\]

 откуда 

\[f= \frac{d F_1}{d + F_1} = \frac{10 \cdot 10}{10 + 10} =5 \,\text{см}.\]

 Это мнимое изображение находится перед рассеивающей линзой. 
$2$.  Учитывая, что данная точка является источником для собирающей линзы, а лучи после прохождения через нее параллельны, делаем вывод, что эта точка находится в фокусе собирающей линзы, то есть на расстоянии $15\,\text{см}$.
$3$.  Тогда расстояние между линзами равно $x=F_2 - f=15 - 5=10 \,\text{см}.$ 

Ответ: $10\,\text{см}$.
}

\end{document}
